\chapter{Task 1 --- CNN Implementation and Model Documentation}
\label{ch:task1}

%-------------------------------------------------------

Task 1 asks for the CNN implementation in PyTorch, the purpose and
input/output dimensions of each layer, a table of trainable parameters
with an explanation of how they were counted, and the training
configuration with a short justification. Section \ref{sec:model} covers
the layers, Section \ref{sec:params} the parameters, Section
\ref{sec:config} the configuration, and Section \ref{sec:loop} the
training loop and how the metrics are calculated.


%%%%%%%%%%%%%%%%%%%%%%%%%%%%%%%%%%%%%%%%%%%%%%%%%%%%%
\section{Architecture and Layer Dimensions}
\label{sec:model}
%%%%%%%%%%%%%%%%%%%%%%%%%%%%%%%%%%%%%%%%%%%%%%%%%%%%%

The network follows the LeNet-5 design \citep{lecun1998}: two
convolutional layers, each followed by ReLU and $2 \times 2$ max pooling,
and then three fully connected layers (Listing \ref{lst:model}). This
meets the minimum requirement of two convolutional layers, one pooling
layer and one fully connected layer.

\lstinputlisting[language=Python, firstline=7, lastline=34,
    firstnumber=auto,
    caption={The CNN (\texttt{model.py}).},
    label={lst:model}]{../model.py}

The purpose of each layer, for one input image, is:

\begin{itemize}
  \item \textbf{conv1} ($1 \times 28 \times 28 \rightarrow 6 \times 28 \times 28$):
        six $5 \times 5$ kernels learn low-level features such as stroke
        edges and line ends, the same kind of features the Sobel filter
        found in Assignment 2. Padding of 2 keeps the size at $28 \times 28$
        (output size $= 28 + 2\cdot2 - 5 + 1 = 28$), so strokes near the
        border are not lost. The lecture example feeds a
        $1 \times 32 \times 32$ image to conv1 without padding
        \citep{abidl08}; padding the $28 \times 28$ MNIST image by 2 pixels
        gives the same $32 \times 32$ input, so every later layer keeps the
        lecture's dimensions.
  \item \textbf{pool1} ($6 \times 28 \times 28 \rightarrow 6 \times 14 \times 14$):
        max pooling halves the resolution and keeps the strongest response
        in every $2 \times 2$ block.
  \item \textbf{conv2} ($6 \times 14 \times 14 \rightarrow 16 \times 10 \times 10$):
        sixteen $5 \times 5$ kernels, each looking at all six maps from
        conv1, combine edges into larger parts of digits such as curves,
        corners and loops. No padding, so $14 - 5 + 1 = 10$.
  \item \textbf{pool2} ($16 \times 10 \times 10 \rightarrow 16 \times 5 \times 5$):
        halves the resolution again.
  \item \textbf{flatten} ($16 \times 5 \times 5 \rightarrow 400$): turns the
        feature maps into one vector for the fully connected layers.
  \item \textbf{fc1} ($400 \rightarrow 120$) and \textbf{fc2}
        ($120 \rightarrow 84$): combine the features from the whole image
        to decide which digit it is.
  \item \textbf{fc3} ($84 \rightarrow 10$): one score per digit. There is
        no softmax here because \texttt{nn.CrossEntropyLoss} applies
        log-softmax internally (Equation \ref{eq:ce}).
\end{itemize}

ReLU is applied after conv1, conv2, fc1 and fc2. ReLU, pooling and
flatten have no parameters.


%%%%%%%%%%%%%%%%%%%%%%%%%%%%%%%%%%%%%%%%%%%%%%%%%%%%%
\section{Trainable Parameters}
\label{sec:params}
%%%%%%%%%%%%%%%%%%%%%%%%%%%%%%%%%%%%%%%%%%%%%%%%%%%%%

A convolutional layer with $C_{in}$ input channels, $C_{out}$ kernels of
size $k \times k$ has $C_{out}(C_{in} k^2 + 1)$ parameters: each kernel
has $C_{in} k^2$ weights plus one bias. A fully connected layer with
$n_{in}$ inputs and $n_{out}$ outputs has $n_{out}(n_{in} + 1)$
parameters. Table \ref{table:params} applies these formulas to every
layer.

\begin{table}[h]
    \centering
    \caption[Trainable parameters per layer.]%
    {Output shape and trainable parameters of every layer, for one input image.}

    \begin{tabular}{c|c|P{2.3in}|r}
        \specialrule{1.5pt}{0pt}{0pt}
        Layer & Output shape & Calculation & Parameters \\
        \hline\hline
        conv1 & $6 \times 28 \times 28$ & $6 \times (1 \cdot 5 \cdot 5 + 1)$ & 156 \\
        \specialrule{0.5pt}{0pt}{0pt}
        pool1 & $6 \times 14 \times 14$ & --- & 0 \\
        \specialrule{0.5pt}{0pt}{0pt}
        conv2 & $16 \times 10 \times 10$ & $16 \times (6 \cdot 5 \cdot 5 + 1)$ & 2,416 \\
        \specialrule{0.5pt}{0pt}{0pt}
        pool2 & $16 \times 5 \times 5$ & --- & 0 \\
        \specialrule{0.5pt}{0pt}{0pt}
        flatten & 400 & --- & 0 \\
        \specialrule{0.5pt}{0pt}{0pt}
        fc1 & 120 & $120 \times (400 + 1)$ & 48,120 \\
        \specialrule{0.5pt}{0pt}{0pt}
        fc2 & 84 & $84 \times (120 + 1)$ & 10,164 \\
        \specialrule{0.5pt}{0pt}{0pt}
        fc3 & 10 & $10 \times (84 + 1)$ & 850 \\
        \hline\hline
        \textbf{Total} & & & \textbf{61,706} \\
        \specialrule{1.5pt}{0pt}{0pt}
    \end{tabular}
    \label{table:params}
\end{table}

To check the hand calculation, \texttt{model.py} also counts the
parameters in code: \texttt{layer\_table} passes a dummy image through
the network with forward hooks to record each layer's output shape, and
sums \texttt{p.numel()} over each layer's tensors. The program prints the
same 61,706 total (Listing \ref{lst:output-params}). Almost 78\% of the
parameters are in fc1, while the two convolutional layers together hold
only about 4\%, because a convolution reuses the same small kernel at
every position of the image.

\begin{lstlisting}[language={}, numbers=none, caption={Program output: layer shapes and parameter counts.}, label={lst:output-params}]
=== Layers (input -> output shape, trainable parameters)
conv1  (1, 28, 28)     -> (6, 28, 28)         156
conv2  (6, 14, 14)     -> (16, 10, 10)       2416
fc1    (400,)          -> (120,)            48120
fc2    (120,)          -> (84,)             10164
fc3    (84,)           -> (10,)               850
total trainable parameters: 61706
\end{lstlisting}


%%%%%%%%%%%%%%%%%%%%%%%%%%%%%%%%%%%%%%%%%%%%%%%%%%%%%
\section{Training Configuration}
\label{sec:config}
%%%%%%%%%%%%%%%%%%%%%%%%%%%%%%%%%%%%%%%%%%%%%%%%%%%%%

\begin{table}[h]
    \centering
    \caption[Baseline training configuration.]%
    {Baseline training configuration.}

    \begin{tabular}{P{1.7in}|P{3.6in}}
        \specialrule{1.5pt}{0pt}{0pt}
        Setting & Value \\
        \hline\hline
        Optimizer & SGD with momentum 0.9 \\
        \specialrule{0.5pt}{0pt}{0pt}
        Learning rate & 0.01 (constant, no scheduler) \\
        \specialrule{0.5pt}{0pt}{0pt}
        Batch size & 64 (750 batches per epoch) \\
        \specialrule{0.5pt}{0pt}{0pt}
        Epochs & 10 \\
        \specialrule{0.5pt}{0pt}{0pt}
        Loss function & Cross-entropy (\texttt{nn.CrossEntropyLoss}) \\
        \specialrule{0.5pt}{0pt}{0pt}
        Weight initialization & PyTorch default, drawn right after seeding \\
        \specialrule{0.5pt}{0pt}{0pt}
        Random seed & 42 (split, initial weights, shuffling) \\
        \specialrule{0.5pt}{0pt}{0pt}
        Preprocessing & ToTensor, then Normalize(mean 0.1307, std 0.3082) \\
        \specialrule{0.5pt}{0pt}{0pt}
        Regularization & None (no dropout, no weight decay) \\
        \specialrule{1.5pt}{0pt}{0pt}
    \end{tabular}
    \label{table:config}
\end{table}

\textbf{Justification.} Cross-entropy is the standard loss for
single-label classification with 10 classes, and it is the loss used in
the lecture's full training example. That example trains with
\texttt{optim.SGD(lr=0.001, momentum=0.9)} and a batch size of 4
\citep{abidl08}. I kept SGD with momentum, which averages recent
gradients and so smooths the noisy mini-batch updates. I used a batch
of 64, which makes each gradient estimate less noisy and an epoch much
faster (750 updates instead of 12,000). With 16 times fewer updates per
epoch, a learning rate of 0.001 would move slowly, so I raised it to
0.01 for the baseline, the value the lecture uses in its
parameter-update example \citep{abidl08}. Task 4 (Chapter
\ref{ch:task4}) tests whether this was a good choice. I left out dropout and weight decay so that the baseline shows
the plain learning behavior of the network, and any overfitting would be
visible in the curves.


%%%%%%%%%%%%%%%%%%%%%%%%%%%%%%%%%%%%%%%%%%%%%%%%%%%%%
\section{Training Loop and Metric Calculation}
\label{sec:loop}
%%%%%%%%%%%%%%%%%%%%%%%%%%%%%%%%%%%%%%%%%%%%%%%%%%%%%

Listing \ref{lst:loop} shows the training loop. For every batch it
clears the old gradients, runs the forward pass, computes the loss,
back-propagates and updates the weights. During the first epoch it also
records, for each batch, the loss, the accuracy and the macro-F1 of that
batch's predictions, plus a cumulative macro-F1 over every image seen so
far in the epoch. These are the values plotted in Task 2 (Chapter
\ref{ch:task2}).

\lstinputlisting[language=Python, firstline=87, lastline=120,
    firstnumber=auto,
    caption={Training loop with first-epoch batch logging (\texttt{train.py}).},
    label={lst:loop}]{../train.py}

At the end of every epoch, \texttt{evaluate} (Listing \ref{lst:evaluate})
runs the model in evaluation mode, without gradients, over the whole
training split and the whole validation split. I measure the training
metrics this way instead of averaging the batch values during the epoch,
because the batch values mix early and late weights from the same epoch.
Re-evaluating both sets with the same end-of-epoch weights makes the
training and validation numbers directly comparable, which matters for
the gap analysis in Task 3 (Chapter \ref{ch:task3}).

Every F1-score in this report is macro-F1 (Equation \ref{eq:macrof1}),
computed by \texttt{macro\_f1} with scikit-learn's
\texttt{f1\_score(average="macro")}: the F1-score of each digit is
computed separately and the values are averaged with equal weight over
the digits that appear in the labels or the predictions. On the full
training, validation and test sets that is always all 10 digits.
Macro-precision and macro-recall for Task 5 are computed the
same way with \texttt{precision\_score} and \texttt{recall\_score}.

\lstinputlisting[language=Python, firstline=32, lastline=57,
    firstnumber=auto,
    caption={Macro-F1 and the evaluation function (\texttt{train.py}).},
    label={lst:evaluate}]{../train.py}
